% ============================================================
% introduction_draft.tex
% Copy contents into Chapters/Introduction.tex
% ============================================================

\newpage
\section{Introduction}

Market practitioners have been reading price charts for over a century. The visual patterns that technical analysts study, including head-and-shoulders formations, support and resistance levels, and momentum breakouts, are still widely used in equity trading despite decades of academic scepticism. The question this thesis asks is simple: can a neural network learn to read those charts, and if it can, does the signal it finds in European equities tell us anything useful?

The starting point is the work of Jiang, Kelly, and Xiu (2023), who convert daily OHLCV data into standardised greyscale images and train convolutional neural networks to predict the direction of 20-day returns. Applied to U.S. stocks from 1993 to 2020, their models produce AUC values in the 0.54--0.56 range and Sharpe Ratios that exceed those of standard momentum strategies. The key idea is that once the data is drawn as an image, a CNN can pick up non-linear patterns across several price features at once, which a hand-built factor model cannot.

This approach, however, has almost exclusively been tested on U.S. equities. It is not obvious that the patterns a CNN learns from NYSE and NASDAQ data will appear in the same form in European markets, which differ in market structure, trading conventions, and investor composition. Nor is it obvious that more modern architectures or richer image formats would improve on the original greyscale chart. And the question of whether visual features trained on equities might transfer to cryptocurrency markets, which share the charting format but differ in nearly every other respect, has not been examined at all.

This thesis addresses all four of these open questions. The empirical universe is the Euro Stoxx 50, the 50 largest Eurozone blue-chip stocks, with daily price data from 2000 to early 2026. Three models are evaluated. The first replicates the baseline CNN of Jiang et al. on European data using the same image format and architecture. The second pairs EfficientNet-B2 with a Transformer sequence encoder that processes 21 technical indicators alongside the chart image. The third pairs ConvNeXt V2-Tiny with an iTransformer encoder that processes 25 indicators. A separate comparison experiment isolates the effect of the image encoding by retraining both hybrid models under two different image formats: a five-panel candlestick chart with multiple overlaid indicators (v3) and a multi-scale RGB image that encodes price dynamics at three time horizons simultaneously (v4).

The main contributions of this thesis are the following. First, it provides out-of-sample evidence that CNN chart pattern models do produce a statistically significant predictive signal in European equities, with AUC values above 0.5 across all architectures. Second, it shows that the image encoding format matters more than the choice of neural network architecture: the multi-scale RGB image significantly outperforms the single-scale candlestick chart for both EfficientNet and ConvNeXt, while the two architectures perform at a similar level within each image type. Third, it shows that a model trained entirely on Euro Stoxx 50 data retains meaningful predictive power when applied to Bitcoin and Ethereum without any retraining, with ConvNeXt V2 achieving an AUC of 0.5433 on the crypto test set. Fourth, it provides a direct replication of Jiang et al. on a non-U.S. market with a strictly walk-forward evaluation design, adding to the limited international evidence on chart-based CNN models.

The remainder of this thesis is structured as follows. Section~2 reviews the relevant literature on machine learning in asset pricing, technical analysis, and image-based financial prediction. Section~3 develops the four empirical hypotheses. Section~4 describes the data, image construction, model architectures, training protocol, and evaluation framework. Section~5 presents and interprets the out-of-sample results for all four hypotheses, including the portfolio analysis and temporal stability tests. Section~6 discusses the main limitations. Section~7 concludes.
